\section*{Abstract}
% \begin{itemize}
%     \item Radar is able to provide Unmanned Aeriel Vehicales with the capability to detect and track airborne hazzards over large distances and in all weather conditions. Talk about trade off between size, power, weight and compute. Talk about problems.
%     \item We seek to address these problems through adapting establish radar processing algorithms for UAVs.
%     \item By evaluating the effectiveness of different implementations and algorithms essential to multi target tracking, this report demonstrates how software choises make long range radar practical for small UAV detect and avoid systems.
%     \item Talk about key parts of our pipeline
%     \item Talk about the Experiments
%     \item Talk about findings
%     \item Ultimately, our work provides a reference for engineers and policy makers...
% \end{itemize}

% Radar is a promising technology that enables collision avoidance 
% systems for \gls{uav}s, offering long-range, all-weather detection
% independent of target cooperation. Recent miniaturisation of millimetre-wave
% radar hardware and the availability of capable embedded \gls{gpu} compute have made
% onboard radar processing feasible for small \gls{uav} platforms. However, existing
% work consistently stops short of validating a complete sensing-to-tracking
% pipeline under realistic airborne conditions, has not benchmarked \gls{gpu}-based
% radar implementations under \gls{uav} \gls{swap} constraints,
% and has only evaluated angle estimation algorithms in idealised noise rather
% than in end-to-end pipelines. This thesis addresses these gaps by designing,
% implementing, and evaluating a software-defined monopulse radar tracking
% pipeline on NVIDIA Jetson Xavier embedded \gls{gpu} platforms, systematically
% comparing algorithm choices and low-level implementations across hardware
% configurations. The pipeline achieves real-time multi-target tracking with
% end-to-end latency below 50\,ms. Taylor windows calibrated to the \gls{cfar} 
% noise
% floor are found to minimise false alarms without recall penalty; coherent
% integration raises detection recall from near zero to 0.9 in the worst noise
% regime tested; and CA-\gls{cfar} outperforms competing variants in high-noise
% conditions. \gls{gpu}-based \gls{cfar} is shown to be the preferred implementation on the
% Xavier platform, while the optimal \gls{cpu}-versus-\gls{gpu} crossover differs between
% embedded and desktop configurations, demonstrating that implementation
% selection cannot be made platform-agnostically. Monopulse angle sensing
% is validated as an unbiased bearing and elevation estimator even at
% $-10$\,dB \gls{snr}, and a Kalman-filter-based multi-target tracker maintains
% position errors below 22\,m RMS across all noise regimes tested. A live
% airborne field trial appears to correctly geolocate static ground targets at up 
% to
% 706\,m, confirming physical pipeline validity but fell short of tracking
% an aerial target due to hardware noise artefacts. Together, these results
% provide hardware-aware design guidance for engineers developing compact radar
% \gls{daa} systems and provide an evidence base for aviation policy.
% Radar enables all-weather, long-range collision avoidance for UAVs independent of target cooperation.
% Recent miniaturisation of millimetre-wave radar hardware and the availability of capable 
% embedded \gls{gpu} compute have made onboard radar processing feasible for small 
% \gls{uav} platforms. Detect-and-avoid, however, demands sensitivity to small targets at 
% ranges of hundreds of metres to several kilometres, and existing work has yet to fully 
% address the requirements this imposes. Complete sensing-to-tracking pipelines have not 
% been validated under realistic airborne conditions; \gls{gpu}-based radar 
% implementations have not been benchmarked under \gls{uav} \gls{swap} constraints; and 
% implementation and it is not clear how algorithm conclusions differ between hardware contexts.
% This thesis addresses 
% these gaps by designing, implementing, and evaluating a software-defined monopulse 
% radar tracking pipeline on desktop and NVIDIA Jetson Xavier embedded \gls{gpu} platforms. Algorithm 
% choices and low-level implementations are compared through ablation studies, parameter 
% sweeps over noise, number of targets, frame size, and detector configuration.

% The pipeline achieves real-time multi-target tracking of three simultaneous targets, 
% sustaining frame rates from 20 to 500 frames per second on the embedded platforms 
% depending on integration size, with end-to-end latency below 50\,ms. Taylor windows 
% calibrated to the \gls{cfar} noise floor minimise false alarms; 
% coherent integration raises detection recall from near zero to 0.9 in the worst noise 
% regime tested; and CA-\gls{cfar} outperforms competing variants in high-noise conditions. 
% Notably, the optimal implementation of algorithms depend on hardware, with the 
% \gls{gpu}-based \gls{cfar} being preferred on the Xaviers but not the desktop. 
% The unified-memory Jetson modules outperform the more powerful desktop in terms of 
% memory management, as their shared \gls{cpu}--\gls{gpu} memory removes the host-device transfer that dominates this workload. 
% Monopulse angle sensing is validated as an unbiased bearing and elevation estimator 
% even at -10\,dB \gls{snr} and a Kalman-filter-based multi-target tracker maintains 
% position errors below 22\,m RMS across all noise regimes, comfortably within the 
% separation distances a detect-and-avoid system must resolve. A live airborne field trial geolocated static ground 
% targets at up to 706\,m, but did not track an 
% aerial target owing to a transmitter-receiver coupling artefact that saturated the near 
% field. Together, these results provide hardware-aware design guidance for engineers 
% developing compact radar \gls{daa} systems and an evidence base to inform aviation policy.

Radar enables all-weather, long-range collision avoidance for UAVs independent of target cooperation. Miniaturised millimetre-wave radar hardware and capable embedded \gls{gpu} compute have made onboard radar processing feasible for small \gls{uav} platforms. Detect-and-avoid, however, demands sensitivity to small targets at ranges of hundreds of metres to several kilometres, and existing work has yet to fully address this. Complete sensing-to-tracking pipelines have not been validated under realistic airborne conditions; \gls{gpu}-based implementations have not been benchmarked under \gls{uav} \gls{swap} constraints; and it is unclear how algorithm conclusions differ between hardware contexts. This thesis addresses these gaps by designing, implementing, and evaluating a software-defined monopulse radar tracking pipeline on desktop and NVIDIA Jetson Xavier embedded \gls{gpu} platforms, comparing algorithm and implementation choices through ablation studies and parameter sweeps over noise, target count, frame size, and detector configuration.

The pipeline achieves real-time tracking of three simultaneous targets at 20 to 500 frames per second on the embedded platforms, with end-to-end latency below 50\,ms. Taylor windows calibrated to the \gls{cfar} noise floor minimise false alarms; coherent integration raises detection recall from near zero to 0.9 in the worst noise regime tested; and CA-\gls{cfar} outperforms competing variants in high noise. The optimal implementation depends on hardware: \gls{gpu}-based \gls{cfar} is preferred on the Xaviers but not the desktop, and the unified-memory Jetson modules outperform the more powerful desktop, as their shared \gls{cpu}--\gls{gpu} memory removes the host-device transfer that dominates this workload. Monopulse angle sensing is validated as an unbiased bearing and elevation estimator even at 
-10
-10\,dB \gls{snr}, and a Kalman-filter tracker maintains position errors below 22\,m RMS across all noise regimes, within the separation distances a detect-and-avoid system must resolve. A live airborne field trial geolocated static ground targets at up to 706\,m, but did not track an aerial target owing to a transmitter-receiver coupling artefact that saturated the near field. These results provide hardware-aware design guidance for engineers developing compact radar \gls{daa} systems and an evidence base to inform aviation policy.